\section{Method}
\label{sec:method}

% Drafted from chebfeatdepth.py. Keep in sync with the code.

Given an RGB image $I \in \mathbb{R}^{H\times W\times 3}$, ChebDepth predicts a dense depth map $\hat{D} \in [0,1]^{H\times W}$ in four stages: image-to-graph conversion, Chebyshev graph feature extraction, graph-to-image projection, and depth decoding (Fig.~\ref{fig:overview}).

\begin{figure*}[t]
  \centering
  % \includegraphics[width=\textwidth]{../figures/overview.pdf}
  \todo{Pipeline overview figure: image $\rightarrow$ superpixel graph $\rightarrow$ ChebNet $\rightarrow$ barycentric projection $\rightarrow$ depth head.}
  \caption{Overview of ChebDepth. \todo{Caption.}}
  \label{fig:overview}
\end{figure*}

\subsection{Image-to-Graph Conversion}
The image is converted to CIELAB and over-segmented with SLICO \cite{achanta2012slic} (region size 20, ruler 10, 10 iterations). The centroid of each superpixel, normalized to $[0,1]^2$, becomes a node $v_i$ with position $\mathbf{p}_i=(x_i,y_i)$.
Each node carries a 5-dimensional feature $\mathbf{f}_i=[R_i,G_i,B_i,x_i,y_i]$, where the color is sampled at the pixel nearest to the centroid.
Edges connect each node to its $k=8$ nearest neighbors in position space (a $k$-NN graph on the centroids, not the superpixel adjacency).
A regular-grid node strategy is also implemented and can serve as an ablation.

\subsection{Chebyshev Graph Feature Extractor}
Node features are processed by four Chebyshev graph convolution layers \cite{defferrard2016chebnet} with polynomial order $K=3$ and widths $64\!\rightarrow\!128\!\rightarrow\!256\!\rightarrow\!512$. Each layer is followed by batch normalization, ReLU, and dropout ($p=0.1$). A filter of order $K$ computes
\begin{equation}
  \mathbf{H}^{(\ell+1)} = \sum_{j=0}^{K-1} T_j(\tilde{\mathbf{L}})\,\mathbf{H}^{(\ell)}\,\mathbf{W}_j^{(\ell)},
\end{equation}
where $T_j$ is the Chebyshev polynomial of the first kind and $\tilde{\mathbf{L}}$ is the rescaled graph Laplacian.
An attention pooling head also produces a global descriptor; \todo{state whether the global feature is used downstream, since the depth head currently consumes only the per-node features}.

\subsection{Graph-to-Image Projection}
Node features are mapped to every pixel by barycentric interpolation over the Delaunay triangulation of the node positions. Pixels outside the triangulation take the feature of the nearest node. The result is a $512\times H\times W$ feature map.

\subsection{Depth Head}
A convolutional decoder with channel widths $256\!\rightarrow\!128\!\rightarrow\!64\!\rightarrow\!32$ (3$\times$3 convolutions, batch normalization, ReLU, dropout, bilinear upsampling in the last stages) followed by a 1$\times$1 convolution and a sigmoid outputs $\hat{D}\in[0,1]$.

\subsection{Training Objective}
The loss combines four terms:
\begin{equation}
  \mathcal{L} = \mathcal{L}_{\mathrm{MSE}} + 0.5\,\mathcal{L}_{\mathrm{grad}} + 0.1\,\mathcal{L}_{\mathrm{smooth}} + 0.05\,\mathcal{L}_{\mathrm{graph}}.
\end{equation}
$\mathcal{L}_{\mathrm{grad}}$ is the $\ell_1$ distance between the horizontal and vertical finite-difference magnitudes of $\hat{D}$ and $D$.
$\mathcal{L}_{\mathrm{smooth}}$ penalizes $|\partial \hat{D}|$ weighted by $e^{-|\partial I|}$, so depth discontinuities are cheaper at image edges.
$\mathcal{L}_{\mathrm{graph}}$ is the mean $\ell_2$ distance between features of connected nodes, a first-order smoothness regularizer in the spirit of Dirichlet energy (it is the unsquared norm; \todo{decide whether to describe it as Dirichlet energy or as a total-variation-style penalty}).
Optimization uses AdamW (learning rate $10^{-4}$, weight decay $10^{-5}$) with a reduce-on-plateau schedule (patience 5, factor 0.5).
